\section{High girth shortens terminal geodesics}
\label{sec:four-high-connectivity-moore-paths}

All logarithms in this section have base two.
A normalized $k$-pole is connected, has $k$ vertices of degree two, and has degree three at every other vertex.
The first bound uses no cyclic-connectivity assumption.

\begin{lemma}[Moore balls along a geodesic]
\label{lem:four-pole-moore-geodesic}
Let $N$ be a normalized $k$-pole of order $n$ and girth at least an integer $g\ge3$.
Put
\[
 t=\left\lfloor\frac{g-1}{2}\right\rfloor,
 \qquad B_t=3\cdot2^t-2.
\]
Every shortest path in $N$, of length $\ell$, satisfies
\begin{equation}\label{eq:four-moore-geodesic-exact}
 \ell+1<(2t+1)\left(k+1+\frac{n}{B_t}\right)
 \le g\left(k+1+n2^{-g/2}\right).
\end{equation}
\end{lemma}

\begin{proof}
Write the path as $v_0v_1\cdots v_\ell$.
Choose centers at positions $t+j(2t+1)$ whenever the full corresponding interval lies in $[0,\ell]$.
There are
\[
 q=\left\lfloor\frac{\ell+1}{2t+1}\right\rfloor
\]
centers.
The radius-$t$ vertex balls around these centers are disjoint.
Otherwise two centers would have distance at most $2t$, contradicting the geodesic property.
At most $k$ balls contain a degree-two vertex.

Every remaining ball contains at least $B_t$ vertices.
Indeed, breadth-first paths of length at most $t$ cannot collide, since a collision gives a cycle of length at most $2t<g$.
All vertices at distance below $t$ have degree three.
The breadth-first tree therefore has $1+3(1+2+\cdots+2^{t-1})=B_t$ vertices.
Edges within the last sphere do not affect this count.
Thus $(q-k)_+B_t\le n$, and hence $q\le k+n/B_t$.
The floor defining $q$ gives the first bound.

Since $t\ge1$, one has $B_t\ge2^{t+1}\ge2^{g/2}$.
Also $2t+1\le g$.
These inequalities give the final bound.
\end{proof}

\begin{corollary}[Short cycles through every edge at high girth]
\label{cor:four-high-girth-edge-cycle}
Let $G$ be a connected bridgeless cubic graph of rank $r$ and girth at least $g\ge3$.
Every edge lies on a cycle of length less than
\[
 g\left(3+(2r-2)2^{-g/2}\right).
\]
In particular, for every fixed $\eta>0$, girth at least $(4+\eta)\log\log r$ gives cycles of length $o(r/\log^2r)$ through every edge.
\end{corollary}

\begin{proof}
For an edge $e=uv$, the graph $N=G-e$ is a connected normalized two-pole of order $2r-2$.
Its girth is at least that of $G$.
A shortest terminal path together with $e$ gives the desired cycle.
Apply Lemma~\ref{lem:four-pole-moore-geodesic} with $k=2$.

For the asymptotic assertion, use the integer lower bound $g_0=\lfloor(4+\eta)\log\log r\rfloor$ in that lemma.
Then $g_0=O(\log\log r)$ and $2^{-g_0/2}=O((\log r)^{-2-\eta/2})$.
Both terms in the asserted length bound are $o(r/\log^2r)$.
\end{proof}

\begin{corollary}[Endpoint forest error for high-girth edge deletion]
\label{cor:four-high-girth-deletion-excess}
There is an absolute constant $C$ with the following property.
Let $G$ be connected and bridgeless cubic, of sufficiently large rank $r$.
Write $x=\log r$ and $z=\log x$.
If
\[
 \operatorname{girth}(G)\ge
 \left\lceil4z+2\log z+12\right\rceil,
\]
then, for every edge $e=uv$,
\[
 \frac{P_{uv}(G-e)}{2^{r-1}}
 \le\frac12+C\frac{(\log r)^{11/2}}{\sqrt r}.
\]
The same conclusion holds under girth at least $(4+\eta)\log\log r$ for any fixed $\eta>0$ and sufficiently large $r$.
\end{corollary}

\begin{proof}
Use $g_0=\lceil4z+2\log z+12\rceil$ in Corollary~\ref{cor:four-high-girth-edge-cycle}.
For sufficiently large $z$, one has $g_0\le8z$ and
\[
 2^{-g_0/2}\le\frac{1}{64x^2z}.
\]
Thus every edge $e=uv$ lies on a cycle $W=A\cup\{e\}$ with
\[
 |V(W)|<24z+\frac{r}{4x^2}\le\frac{r}{x^2}.
\]
Here $A$ is a terminal path in $G-e$.
Each terminal path $P$ determines a distinct outside part $P\setminus E(A)$.
Distinctness follows because their symmetric difference cannot be a nonempty even set supported on the path $A$.
There are $2^{r-1}$ possible outside parts, and each path gives a forest outside $W$.

The outside-part law is the restriction of the uniform cycle-space law of $G$.
Being a forest outside the edges of $W$ implies being a forest outside its vertices.
Apply Theorem~\ref{thm:parallel-links-log-five-half-forest} with the connected marked set $V(W)$.
Every vertex of $G$ is cubic, and the displayed size bound supplies its remaining hypothesis.
This proves the estimate uniformly over all edges.
\end{proof}

\begin{proposition}[A rank-independent girth bound for path excess]
\label{prop:four-girth-path-excess}
There is an absolute constant $C$ such that every normalized two-pole $N$ of rank $b$ and girth at least $g\ge3$ satisfies
\[
 \left(\frac{P_{uv}(N)}{2^b}-\frac12\right)_+
 \le C g^{11/2}2^{-g/4}.
\]
\end{proposition}

\begin{proof}
Let $t=\lfloor(g-1)/2\rfloor$, and let $L_t$ be the distance-$t$ sphere from the terminal $u$.
The breadth-first expansion starts with two edges at $u$.
Until radius $t$, distinct expansion paths cannot reach the same vertex, because $2t<g$.
Only the other terminal $v$ can remove an expected child.
If its distance from $u$ is $s<t$, that removal costs $2^{t-s-1}$ vertices in the final sphere.
Otherwise there is no removal.
Consequently
\[
 |L_t|\ge2^{t-1}\ge2^{g/2-2}.
\]

Suppose the normalized path count is at least $1/2+\varepsilon$.
The quantitative localization proof applies directly to $N$, with no cubic expansion needed.
Its distance-layer bound, from Theorem~\ref{thm:path-quantitative-additive}, gives
\[
 |L_t|\le K_1\varepsilon^{-2}\bigl(\log(2/\varepsilon)\bigr)^{11}+1.
\]
Choose $C\ge4\sqrt{K_1+1}$ and put $\varepsilon=Cg^{11/2}2^{-g/4}$.
For sufficiently large $g$, this lies in $(0,1/2]$ and $\log(2/\varepsilon)\le g$.
The last upper bound is then at most
\[
 \frac1{16}2^{g/2}+1<2^{g/2-2},
\]
a contradiction.
Increase $C$ to cover the finitely many remaining values of $g$.
\end{proof}

The independent geometric checker is \texttt{four-high-connectivity-moore-check.py}.
Its output, \texttt{four-high-connectivity-moore-checks.json}, records 49 normalized pole fixtures.
They include edge-deleted and subdivided Petersen, Heawood, and honeycomb graphs, with girth up to twelve.
All 35,568 geodesic checks and 40 terminal-sphere checks pass.
These finite checks validate the counting implementation; the proof above supplies the general statement.
